% Auto-generated answer document. Original problem statements are preserved verbatim.
\documentclass[11pt]{article}
\usepackage{fontspec}
\usepackage{microtype}
\usepackage{mathtools,amssymb,amsfonts,amsthm,bm}
\usepackage{enumitem}
\usepackage{geometry,fancyhdr,setspace,xcolor}
\usepackage{fancyvrb}
\usepackage{pdfpages}
\usepackage[hidelinks]{hyperref}
\geometry{margin=1in,headheight=14pt}
\setstretch{1.15}
\setlength{\parindent}{0pt}
\setlength{\parskip}{0.55em plus 0.12em minus 0.08em}
\setlength{\abovedisplayskip}{0.85em plus 0.2em minus 0.1em}
\setlength{\belowdisplayskip}{0.85em plus 0.2em minus 0.1em}
\setlist{leftmargin=*,topsep=0.35em,itemsep=0.15em,parsep=0pt}
\pagestyle{fancy}
\fancyhf{}
\fancyhead[L]{\coursename}
\fancyhead[C]{\assignment}
\fancyhead[R]{\studentname}
\fancyfoot[C]{\thepage}
\newcommand{\studentname}{Keshav Ramamurthy}
\newcommand{\coursename}{Math 118}
\newcommand{\assignment}{Homework 1}
\newcommand{\term}{Fall 2026}
\newcommand{\R}{\mathbb{R}}
\newcommand{\C}{\mathbb{C}}
\newcommand{\Q}{\mathbb{Q}}
\newcommand{\Z}{\mathbb{Z}}
\newcommand{\N}{\mathbb{N}}
\newcommand{\F}{\mathbb{F}}
\newcommand{\K}{\mathbb{K}}
\newcommand{\E}{\mathbb{E}}
\newcommand{\Prob}{\mathbb{P}}
\newcommand{\abs}[1]{\left\lvert#1\right\rvert}
\newcommand{\norm}[1]{\left\lVert#1\right\rVert}
\newcommand{\ip}[2]{\left\langle#1,#2\right\rangle}
\newcommand{\set}[1]{\left\{#1\right\}}
\newcommand{\given}{\,\middle\vert\,}
\newcommand{\dd}{\,\mathrm{d}}
\DeclareMathOperator{\rank}{rank}
\DeclareMathOperator{\nullity}{nullity}
\DeclareMathOperator{\Span}{span}
\DeclareMathOperator{\Ker}{ker}
\DeclareMathOperator{\im}{im}
\DeclareMathOperator{\Hom}{Hom}
\DeclareMathOperator{\Aut}{Aut}
\DeclareMathOperator{\Var}{Var}
\DeclareMathOperator{\Cov}{Cov}
\newtheorem{problem}{Problem}[section]
\newenvironment{solution}{%
  \par\smallskip\noindent\textbf{Solution notes}\par
  \vspace{1.15in}\noindent\hrulefill\par\vspace{0.28in}%
}{\par}
\newenvironment{originalproblem}{%
  \VerbatimEnvironment
  \par\medskip\noindent\begin{minipage}{0.98\linewidth}
  \begin{Verbatim}[fontsize=\small]%
}{%
  \end{Verbatim}\end{minipage}\par
}
\renewcommand{\coursename}{Math 118}
\renewcommand{\assignment}{Homework 1 Answers}
\begin{document}
\begin{center}\Large\textbf{Math 118 -- Homework 1 Answers}\end{center}
\includepdf[pages=1-2,pagecommand={\thispagestyle{empty}}]{hw01.pdf}
\clearpage
\section*{Solutions}
\subsubsection*{Problem 1}
\begin{solution}
    By the triangle inequality,
    \[\norm{u} = \norm{(u - v) + v} \le \norm{u - v} + \norm{v}\]
    and, since $\norm{v - u} = \norm{u - v}$ by absolute homogeneity,
    \[\norm{v} = \norm{(v - u) + u} \le \norm{u - v} + \norm{u}.\]
    Rearranging these two gives
    \[-\norm{u - v} \le \norm{u} - \norm{v} \le \norm{u - v},\]
    which is what we wanted.
% Write your proof, calculation, or argument here.
\end{solution} 
\subsubsection*{Problem 2}
\begin{solution}
    Since commutativity holds over addition and multiplication,
    \[\langle u, v\rangle = \sum_{i=1}^n u_i \overline{v_i}
    = \overline{\sum_{i=1}^n v_i \overline{u_i}} = \overline{\langle v, u\rangle}.\]
    Linearity on the first term and on the second term:
    \begin{align}
    \langle au+bw,v\rangle
    &=\sum_{i=1}^n (au_i+bw_i)\overline{v_i}
    =a\langle u,v\rangle+b\langle w,v\rangle,\\
    \langle u,av+bw\rangle
    &=\sum_{i=1}^n u_i\left(\overline{a}\,\overline{v_i}+\overline{b}\,\overline{w_i}\right)
    =\overline{a}\langle u,v\rangle+\overline{b}\langle u,w\rangle.
    \end{align}
    And for positive definiteness, $\langle u, u\rangle = \sum_{i=1}^n \abs{u_i}^2 \ge 0$, which is 0 only when $u = 0$.

    So this is an inner product on $\C^n$.
% Write your proof, calculation, or argument here.
\end{solution}
\subsubsection*{Problem 3}
\begin{solution}
    With $v_2 = -\frac{1}{2}v_1$,
\begin{align}
\langle V,V\rangle
&=
(\overline{v_1},\overline{v_2})
\begin{pmatrix}
1 & 2\\
2 & 4
\end{pmatrix}
\begin{pmatrix}
v_1\\
v_2
\end{pmatrix}
=
(\overline{v_1},\overline{v_2})
\begin{pmatrix}
v_1+2v_2\\
2v_1+4v_2
\end{pmatrix}\\
&=
(\overline{v_1},\overline{v_2})
\begin{pmatrix}
0\\
0
\end{pmatrix}
=0.
\end{align}
So $\langle V, V\rangle = 0$ whenever $v_1 + 2v_2 = 0$, for example $V = \binom{2}{-1} \ne 0$.
So positive definiteness fails, and this is not an inner product.
% Write your proof, calculation, or argument here.
\end{solution}
\subsubsection*{Problem 4}
\begin{solution}
Compute the null space:
\begin{align}
\mathbf{x}\cdot
\begin{pmatrix}
1\\
-2\\
1
\end{pmatrix}
&=0\\
x-2y+z&=0\\
x&=2y-z
\end{align}

Thus,
\[
\begin{pmatrix}
x\\
y\\
z
\end{pmatrix}
=
y\begin{pmatrix}
2\\
1\\
0
\end{pmatrix}
+
z\begin{pmatrix}
-1\\
0\\
1
\end{pmatrix}.
\]

Therefore,
\[
W^\perp
=
\operatorname{span}
\left\{
\begin{pmatrix}
2\\
1\\
0
\end{pmatrix},
\begin{pmatrix}
-1\\
0\\
1
\end{pmatrix}
\right\}.
\]
\end{solution}
\subsubsection*{Problem 5}
\begin{solution}
Since $f(t) = 1$ on $[0,1]$,
\[\langle f, g\rangle = \int_0^1 g(t)\,dt,\]
which is 0 exactly when $g$ has average value 0 on $[0,1]$ (the average is just $\frac{1}{1-0}$ times this integral).

So $f^\perp$ is the set of functions with average value zero.
% Write your proof, calculation, or argument here.
\end{solution}

\subsubsection*{Problem 6}
\begin{solution}
    Integrating by parts and using $\int_0^\pi f(t)\cos(t)\,dt = 0$,
    \begin{align}
    \langle f',\sin(t)\rangle
    &=\left[f(t)\sin(t)\right]_0^\pi-\int_0^\pi f(t)\cos(t)\,dt\\
    &=0-0=0,
    \end{align}
    since $\sin(0) = \sin(\pi) = 0$.

    So $f'$ is orthogonal to $\sin(t)$ in $L^2[0,\pi]$.
% Write your proof, calculation, or argument here.
\end{solution}

\subsubsection*{Problem 7}
\begin{solution}
Suppose $\ip{u_0}{v} = \ip{u_1}{v}$ for all $v \in V$.
Take $v = u_0 - u_1$ and subtract, using linearity:
\[\ip{u_0 - u_1}{u_0 - u_1}
= \ip{u_0}{u_0 - u_1} - \ip{u_1}{u_0 - u_1} = 0.\]
By positive definiteness this only happens when $u_0 - u_1 = 0$, so $\boxed{u_0 = u_1}$.
\end{solution}
\subsubsection*{Problem 8}
\begin{solution}
Let $\set{u_1, \dots, u_n}$ be orthonormal, so $\ip{u_i}{u_j} = 0$ for $i \ne j$ and $\ip{u_j}{u_j} = 1$.
If $c_1 u_1 + \cdots + c_n u_n = 0$, take the inner product with $u_j$ and use linearity:
\[0 = \sum_{i=1}^n c_i\, \ip{u_i}{u_j} = c_j.\]
So every $c_j = 0$ and the vectors are linearly independent.
An infinite set counts as linearly independent when every finite subset is, and any finite subset of the orthonormal set is orthonormal, so it is linearly independent.
\end{solution}
\end{document}
